\section{Experiencia internacional en la regulación del desempeño de los prestadores}
\label{anexo:experiencia_internacional}

Este anexo examina cinco experiencias regulatorias para identificar qué resultados se exigen o evalúan respecto de los prestadores, cómo se distinguen de las inversiones, actividades, capacidades de gestión y variables financieras utilizadas para producirlos, y qué consecuencias se asocian a su incumplimiento. La comparación no busca establecer equivalencias institucionales exactas con el régimen peruano, sino determinar si los sistemas examinados integran objetos regulatorios heterogéneos dentro de una misma medida o asignan instrumentos diferenciados a la evaluación del desempeño, la verificación de inversiones y actividades y la aplicación de consecuencias económicas o sancionadoras.

Los casos difieren tanto en la referencia utilizada para evaluar el desempeño como en las consecuencias del incumplimiento. Italia emplea clases y objetivos diferenciados; Inglaterra y Gales fija compromisos específicos dentro de cada revisión tarifaria; Portugal recurre a una evaluación multidimensional comparativa; Chile utiliza un ranking normalizado de desempeño relativo; y Colombia establece estándares generales acompañados de metas anuales y gradualidades. Algunos de estos esquemas incorporan penalidades o ajustes económicos vinculados con resultados específicos; otros separan la evaluación comparativa de los procedimientos destinados a establecer infracciones. Inglaterra y Gales ofrece el referente más próximo para examinar la separación entre resultados e inversiones, mientras que Italia y Portugal aportan, respectivamente, antecedentes sobre la heterogeneidad interna de los indicadores y el uso de reconocimientos reputacionales separados de la evaluación ordinaria.

\subsection{Italia: calidad técnica, objetivos diferenciados e incentivos por desempeño}
\label{anexo:italia}

La \emph{Autorità di Regolazione per Energia Reti e Ambiente} (ARERA) define la metodología tarifaria, los objetivos de calidad técnica y contractual y los mecanismos de premios y penalidades del servicio hídrico integrado. El periodo 2024--2029 se rige por el \emph{Metodo Tariffario Idrico} MTI-4 y por la \emph{Regolazione della Qualità Tecnica del Servizio Idrico Integrato} (RQTI). Las autoridades territoriales determinan la posición inicial de cada operador, incorporan los objetivos en la programación y presentan las propuestas tarifarias y de inversión \citep{ARERA2023MTI4,ARERA2023RQTI,ARERA2025Assetto}.

La RQTI combina prerrequisitos de información, estándares específicos aplicables a prestaciones individuales y estándares generales representados mediante siete macroindicadores. El valor observado determina una clase inicial y un objetivo de mantenimiento o mejora; los resultados se reportan anualmente y los premios o penalidades se calculan sobre el desempeño acumulado de cada bienio \citep{ARERA2017RQTI,ARERA2023RQTI}.

\begin{table*}[!t]
\centering
\caption{Macroindicadores de calidad técnica utilizados en Italia}
\label{tab:italia_desagregacion_macroindicadores}
\scriptsize
\renewcommand{\arraystretch}{1.20}
\setlength{\tabcolsep}{5pt}
\begin{tabularx}{\textwidth}{
>{\raggedright\arraybackslash}p{2.65cm}
>{\raggedright\arraybackslash}p{5.55cm}
>{\raggedright\arraybackslash}X}
\hline
\textbf{Macroindicador} &
\textbf{Componentes principales} &
\textbf{Objeto directamente observado} \\
\hline

\textbf{M0. Resiliencia hídrica}
&
Relación entre las necesidades totales y la disponibilidad hídrica de la gestión y del ámbito territorial superior.
&
Suficiencia de las fuentes y capacidad para atender las necesidades previstas.
\\

\textbf{M1. Pérdidas de agua}
&
Pérdidas lineales por kilómetro de red y pérdidas porcentuales respecto del volumen ingresado.
&
Agua perdida en relación con la red y con el volumen introducido al sistema.
\\

\textbf{M2. Interrupciones}
&
Duración de interrupciones programadas y no programadas, ponderada por usuarios afectados.
&
Tiempo medio de afectación del abastecimiento.
\\

\textbf{M3. Calidad del agua}
&
Órdenes de no potabilidad, muestras no conformes y parámetros con resultados no conformes.
&
Incidencia de restricciones al consumo e incumplimientos sanitarios.
\\

\textbf{M4. Adecuación del alcantarillado}
&
Inundaciones o vertimientos, adecuación de aliviaderos e inspección o detección de su activación.
&
Incidentes, condición de la infraestructura y actividades de inspección y control.
\\

\textbf{M5. Gestión de lodos}
&
Proporción de lodos de depuración destinada a relleno sanitario.
&
Destino de los lodos producidos.
\\

\textbf{M6. Calidad del efluente}
&
Muestras de efluentes que superan los límites aplicables.
&
Incumplimientos de las descargas de las plantas de tratamiento.
\\
\hline
\end{tabularx}

\begin{minipage}{\textwidth}
\footnotesize
\textit{Nota:} La tabla presenta los principales componentes utilizados para calcular y clasificar los macroindicadores. La RQTI contempla además indicadores complementarios y prerrequisitos relacionados con la disponibilidad, calidad y confiabilidad de la información. \textit{Fuente:} Elaboración propia sobre la base de la RQTI \citep{ARERA2017RQTI,ARERA2023RQTI}.
\end{minipage}
\end{table*}

M1, M2, M3, M5 y M6 observan principalmente pérdidas, continuidad, conformidad sanitaria o ambiental y gestión de subproductos. M0 incorpora tanto condiciones bajo control del operador como restricciones territoriales de disponibilidad del recurso. M4 es el macroindicador más heterogéneo, pues reúne incidentes del alcantarillado, adecuación de aliviaderos y actividades de inspección o detección. La separación entre resultados y medios no es, por tanto, completa en todos los casos, aunque sus componentes permanecen identificados dentro del sistema \citep{ARERA2017RQTI,ARERA2023RQTI}.

Los objetivos dependen de la clase inicial. Los operadores situados en la clase superior deben mantener el nivel alcanzado, mientras que los demás siguen trayectorias de mejora diferenciadas. Italia utiliza así umbrales nacionales comunes para clasificar el desempeño, sin imponer una única meta absoluta a todos los prestadores \citep{ARERA2017RQTI,ARERA2023RQTI}.

El mecanismo de incentivos conserva esa desagregación: el cumplimiento se evalúa por macroindicador y puede dar lugar a premios o penalidades. Los niveles avanzado y de excelencia incorporan, además, comparaciones sectoriales basadas en el nivel alcanzado y en la mejora registrada. De este modo, el incumplimiento de una trayectoria de mantenimiento o mejora puede generar una consecuencia económica asociada al macroindicador correspondiente. Aunque sus efectos pueden trasladarse posteriormente a las determinaciones tarifarias, cada dimensión mantiene una medición identificable y no se reduce a un promedio único con compensación irrestricta \citep{ARERA2023RQTI,ARERA2023MTI4}.

La calidad se articula con el programa de intervenciones, pero no sustituye su control. El MTI-4 reconoce costos de actuaciones aprobadas, verifica gastos y obras en curso y, para 2024--2029, prevé una penalidad autónoma por incumplimiento de la planificación de inversiones, diferenciada de las consecuencias aplicadas por calidad técnica o contractual \citep{ARERA2023MTI4}. La RQTI evalúa así el desempeño por dimensiones y asocia consecuencias económicas a su evolución, mientras que el MTI-4 controla por separado la programación y ejecución de las inversiones.

\subsection{Inglaterra y Gales: compromisos de desempeño y entregables de inversión}
\label{anexo:ofwat}

La \emph{Water Services Regulation Authority} (Ofwat) establece cada cinco años los ingresos permitidos, las asignaciones de gasto, los compromisos de desempeño y los mecanismos de incentivo de las empresas de agua y saneamiento. El periodo vigente corresponde al \emph{Price Review 2024} (PR24), aplicable entre abril de 2025 y marzo de 2030. La supervisión sanitaria y ambiental corresponde a autoridades especializadas, mientras que Ofwat concentra la regulación económica \citep{OfwatIndustryOverview,Ofwat2024Approach}.

El marco de resultados utiliza \emph{performance commitments}. Cada compromiso cuenta con una definición, unidad de medida, regla de cálculo, nivel comprometido y, cuando corresponde, una tasa de incentivo financiero. PR24 comprende compromisos comunes y medidas específicas para determinadas empresas; en conjunto, cubren experiencia del usuario, calidad y continuidad, funcionamiento del alcantarillado, eficiencia hídrica, ambiente y condición de los activos \citep{Ofwat2024Outcomes,Ofwat2025Definitions}.

\begin{table*}[!t]
\centering
\caption{Compromisos comunes de desempeño establecidos por Ofwat para 2025--2030}
\label{tab:ofwat_compromisos_pr24}
\scriptsize
\renewcommand{\arraystretch}{1.18}
\setlength{\tabcolsep}{5pt}
\begin{tabularx}{\textwidth}{
>{\raggedright\arraybackslash}p{3.05cm}
>{\raggedright\arraybackslash}p{5.65cm}
>{\raggedright\arraybackslash}X}
\hline
\textbf{Dimensión} &
\textbf{Compromisos representativos} &
\textbf{Objeto directamente observado} \\
\hline

\textbf{Usuarios y participantes del mercado}
&
C-MeX, D-MeX, BR-MeX y experiencia de clientes empresariales.
&
Experiencia de usuarios residenciales, urbanizadores, comercializadores y clientes no residenciales.
\\

\textbf{Calidad y continuidad}
&
Interrupciones del suministro, riesgo de cumplimiento sanitario y contactos sobre calidad del agua.
&
Duración de interrupciones, riesgo sanitario y problemas percibidos en el agua suministrada.
\\

\textbf{Alcantarillado}
&
Inundaciones internas y externas.
&
Incidentes que afectan inmuebles, terrenos y espacios exteriores.
\\

\textbf{Demanda y eficiencia hídrica}
&
Fugas, consumo residencial por habitante y demanda de usuarios empresariales.
&
Pérdidas en la red y utilización del recurso por categorías de usuarios.
\\

\textbf{Resultados ambientales}
&
Biodiversidad; emisiones; contaminación; permisos de descarga; aguas de baño y ríos; aliviaderos de tormenta.
&
Resultados ambientales, incidentes y cumplimiento de obligaciones sobre cuerpos receptores.
\\

\textbf{Condición de los activos}
&
Roturas de tuberías, indisponibilidad no programada y colapsos de alcantarillado.
&
Fallas o condiciones de los activos que afectan la confiabilidad del servicio.
\\
\hline
\end{tabularx}

\begin{minipage}{\textwidth}
\footnotesize
\textit{Nota:} Los compromisos se aplican según el tipo de servicio prestado. La agrupación resume las medidas comunes de PR24 y no reproduce todas sus reglas de cálculo. \textit{Fuente:} Elaboración propia sobre la base de las determinaciones finales y definiciones de PR24 \citep{Ofwat2024Outcomes,Ofwat2025Definitions}.
\end{minipage}
\end{table*}

Los niveles comprometidos de desempeño o \emph{performance commitment levels} pueden ser comunes o específicos por empresa y adoptar un valor único para el quinquenio o una trayectoria anual. Su fijación considera el desempeño histórico, la frontera sectorial, las obligaciones legales y ambientales y el gasto eficiente reconocido. Una misma dimensión puede, por ello, traducirse en exigencias distintas según las condiciones y el punto de partida de cada operador \citep{Ofwat2024Outcomes}.

Los \emph{outcome delivery incentives} (ODI) vinculan la desviación respecto del nivel comprometido con pagos o descuentos calculados separadamente para cada compromiso. Ofwat puede establecer incentivos simétricos, mecanismos aplicables solo al desempeño inferior y límites como \emph{caps}, \emph{collars} o \emph{deadbands}. Así, el incumplimiento de un compromiso de calidad o continuidad puede producir una consecuencia financiera específica sin depender del cumplimiento agregado de los demás compromisos. La arquitectura conserva la identificación de cada resultado antes de agregar sus efectos económicos \citep{Ofwat2024Outcomes}.

Los \emph{price control deliverables} (PCD) protegen a los usuarios cuando una asignación destinada a mejoras no queda suficientemente cubierta por un compromiso de desempeño y su ODI. Cada PCD puede especificar el producto financiado, la cantidad comprometida, el perfil temporal, la fecha de culminación y la evidencia requerida. Cuando el beneficio no puede observarse directamente durante el periodo, la verificación recurre a capacidad instalada, infraestructura puesta en operación, población beneficiada, actuaciones completadas o equipos habilitados \citep{Ofwat2023PCD,Ofwat2025PCD}.

El incumplimiento total o parcial permite reconciliar o devolver el financiamiento correspondiente a la parte no entregada; los retrasos pueden generar, además, ajustes por el periodo durante el cual los usuarios no recibieron el beneficio. Ofwat examina expresamente la posible superposición entre PCD, ODI, mecanismos de reparto de costos y obligaciones sanitarias o ambientales para evitar consecuencias duplicadas por un mismo incumplimiento \citep{Ofwat2023PCD,Ofwat2025PCD}.

Los ODI traducen, por tanto, la desviación de los resultados comprometidos en consecuencias financieras, mientras que los PCD controlan productos, hitos, cantidades y puesta en operación de las inversiones financiadas. Ambos se determinan dentro de la misma revisión tarifaria, pero mantienen objetos de evaluación y consecuencias diferenciados.

\subsection{Portugal: evaluación multidimensional y regulación por comparación}
\label{anexo:portugal}

La \emph{Entidade Reguladora dos Serviços de Águas e Resíduos} (ERSAR) regula la calidad, la sostenibilidad económica y la protección de los usuarios de los servicios de agua y residuos. La prestación comprende sistemas mayoristas y minoristas administrados por municipios, entidades intermunicipales, empresas públicas o concesionarios. La evaluación anual de calidad se mantiene diferenciada de la aprobación o supervisión tarifaria correspondiente a cada modalidad de gestión \citep{ERSAR2025RASARP,ERSAR2026Mision,ERSAR2026Tarifas}.

El sistema portugués opera principalmente como una evaluación multidimensional y comparativa. Para cada entidad se publica el valor observado, una referencia y la calificación correspondiente. Estos elementos permiten comparar prestadores y seguir su evolución, pero no constituyen por sí mismos una meta nacional uniforme. La cuarta generación comprende acceso y atención a los usuarios, calidad sanitaria y ambiental, pérdidas y eficiencia operativa, rehabilitación y gestión de activos, cobertura de gastos, circularidad, resiliencia y calidad de la información \citep{ERSAR2021Guia27,ERSAR2026Qualidade}.

La evaluación combina resultados sanitarios, ambientales y operativos con información sobre rehabilitación, gestión de activos, cobertura de gastos y capacidades de planificación. Cada objeto permanece identificado mediante indicadores y calificaciones separadas, en lugar de integrarse necesariamente en un único promedio compensatorio.

Los resultados se publican mediante fichas por entidad. La difusión se complementa con los \emph{Prémios e Selos de Qualidade dos Serviços de Águas e Resíduos}, que reconocen desempeños destacados en categorías como agua segura, pérdidas, continuidad, respuesta a reclamos y cobertura de gastos. Estos premios operan al margen de la calificación ordinaria y constituyen un mecanismo reputacional separado y acotado \citep{ERSAR2026Pesquisa,ERSAR2026Premios}.

Una calificación comparativa desfavorable no constituye, por sí sola, una infracción ni activa automáticamente una sanción. ERSAR dispone, sin embargo, de potestades de regulación, supervisión y fiscalización sobre las obligaciones aplicables a los prestadores. Cuando un resultado deficiente coincide con el incumplimiento de un estándar u obligación vinculante, la consecuencia jurídica deriva de esa obligación y del procedimiento correspondiente, no de la posición obtenida en la evaluación anual \citep{ERSAR2026Mision,ERSAR2026Qualidade}.

La información sobre inversiones, rehabilitación y activos se utiliza en la evaluación económica y en los instrumentos aplicables a cada modalidad de gestión. Las obligaciones concretas de inversión dependen, sin embargo, de contratos, planes, decisiones tarifarias o mandatos municipales y no se sustituyen por la calificación comparativa. Portugal ofrece así un antecedente para diferenciar indicadores económico-financieros, variables de diagnóstico, obligaciones exigibles y mecanismos reputacionales de reconocimiento.

\subsection{Chile: indicadores de calidad y control de los planes de desarrollo}
\label{anexo:chile}

La Superintendencia de Servicios Sanitarios (SISS) regula y fiscaliza a las concesionarias urbanas, interviene en la determinación tarifaria y revisa los planes de desarrollo. Las tarifas se fijan por periodos de cinco años mediante el modelo de una empresa eficiente. Por otra parte, cada concesionaria debe mantener un plan de desarrollo con horizonte de quince años, actualizado integralmente cada cinco, que identifica obras, cronogramas y necesidades de inversión \citep{ChileDFL70,ChileDFL382,SISS2016GuiaPlanes,SISSModeloTarifario}.

La SISS publica anualmente un sistema normalizado de indicadores de calidad. Los valores próximos a uno representan un mejor desempeño relativo, pero la posición en el ranking no equivale automáticamente a la existencia o inexistencia de una infracción: el cumplimiento normativo se determina mediante estándares y procedimientos de fiscalización separados. La evaluación comprende calidad del agua, continuidad y presión del abastecimiento, funcionamiento del alcantarillado, tratamiento de aguas servidas, exactitud del cobro, reclamos y condición de determinados equipos \citep{SISSIndicadoresCalidad,SISS2025Sector}.

La comparación reúne resultados sanitarios, operativos, ambientales, comerciales y de atención, y permite identificar tendencias y diferencias relativas entre concesionarias. Un deterioro de la posición en el ranking no activa por sí mismo una sanción. Si el resultado observado supone, en cambio, el incumplimiento de una obligación exigible de calidad, continuidad u otra condición de prestación, la SISS puede aplicar los procedimientos de fiscalización y sanción previstos en el marco sectorial. En paralelo, revisa los planes de desarrollo y verifica el cumplimiento de las obras y plazos comprometidos \citep{ChileDFL382,SISSIndicadoresCalidad,SISS2025Sector}.

Las inversiones y sus montos se reportan como información económica y operacional, pero no forman parte del ranking anual. Chile mantiene así separados el índice comparativo de calidad, el control normativo de los resultados y la fiscalización de las obras incluidas en los planes de desarrollo.

\subsection{Colombia: metas tarifarias, estándares de prestación y buenas prácticas}
\label{anexo:colombia}

La Comisión de Regulación de Agua Potable y Saneamiento Básico (CRA) establece las metodologías tarifarias, mientras que la Superintendencia de Servicios Públicos Domiciliarios ejerce la inspección, vigilancia y control. La Resolución CRA N.\textdegree~1032 de 2026 introdujo una nueva metodología para prestadores urbanos de más de 5\,000 suscriptores. El régimen comenzó a aplicarse en julio de 2026 y distingue segmentos y condiciones especiales que inciden en las gradualidades y los incentivos \citep{ColombiaLey142,CRA2026Marco}.

Cada prestador debe fijar metas anuales para alcanzar o mantener los estándares de la CRA, incorporarlas en sus estudios de costos y tarifas, reportarlas al SUI y reflejarlas en el contrato de condiciones uniformes. Los estándares generales incluyen, entre otros, 100\,\% de cobertura y 24 horas diarias de continuidad; sin embargo, su cumplimiento se organiza mediante trayectorias anuales y gradualidades según el segmento, el área y las condiciones especiales del prestador. El marco comprende además calidad del agua, cobertura de alcantarillado, calidad del vertimiento, tratamiento y satisfacción del usuario \citep{CRA2026Marco}.

La metodología combina estos resultados con planes e instrumentos de gestión, como la proyección de demanda, la reducción de pérdidas, la sostenibilidad hídrica, la gestión social, la energía, el reúso y la innovación. Las buenas prácticas pueden ser objeto de publicación o de reconocimientos optativos. Para determinados prestadores, el incumplimiento de metas genera descuentos anuales en los componentes de administración, operación o inversión, calculados a partir de los resultados del año anterior \citep{CRA2026Marco}. Para esos indicadores, por tanto, el incumplimiento tiene una consecuencia tarifaria propia.

El Plan de Obras e Inversiones Regulado identifica los activos programados y su año de entrada en operación. Cada año se compara la inversión planeada con la ejecutada y se ajusta el saldo pendiente por recuperar. Este control se mantiene diferenciado de los descuentos asociados con las metas de cobertura, continuidad, calidad o satisfacción. Colombia combina así metas anuales con gradualidad, consecuencias económicas por resultados y un mecanismo separado para verificar los activos planeados, ejecutados y puestos en operación.

\begin{table*}[!t]
\centering
\caption{Síntesis de las arquitecturas internacionales examinadas}
\label{tab:anexo_experiencia_internacional}
\scriptsize
\renewcommand{\arraystretch}{1.16}
\setlength{\tabcolsep}{3.5pt}
\begin{tabularx}{\textwidth}{
>{\raggedright\arraybackslash}p{1.68cm}
>{\raggedright\arraybackslash}p{2.47cm}
>{\raggedright\arraybackslash}X
>{\raggedright\arraybackslash}p{3.20cm}
>{\raggedright\arraybackslash}p{2.85cm}}
\hline
\textbf{Jurisdicción} &
\textbf{Lógica de fijación o evaluación} &
\textbf{Variables comprendidas} &
\textbf{Consecuencia del incumplimiento del desempeño} &
\textbf{Tratamiento de inversiones} \\
\hline

Italia
&
Clases y objetivos de mantenimiento o mejora según la posición inicial.
&
Resiliencia, pérdidas, interrupciones, calidad, alcantarillado, lodos y efluentes.
&
Penalidades económicas por macroindicador, según la clase y la trayectoria aplicable.
&
Programación tarifaria y penalidad autónoma por incumplimiento de la planificación.
\\

Inglaterra y Gales
&
Compromisos comunes o específicos por empresa, con niveles y trayectorias verificables.
&
Usuarios, continuidad, fugas, calidad, ambiente y condición de activos.
&
Los ODI generan pagos o descuentos según la desviación respecto del compromiso; la consecuencia se calcula por compromiso.
&
Los PCD verifican productos, cantidades, hitos, fechas y puesta en operación.
\\

Portugal
&
Valores de referencia y calificaciones comparativas por entidad.
&
Acceso, calidad, pérdidas, activos, cobertura de gastos, circularidad y resiliencia.
&
La calificación comparativa no activa por sí sola una sanción; las obligaciones vinculantes se controlan separadamente.
&
Las obligaciones concretas dependen de planes, contratos y decisiones tarifarias.
\\

Chile
&
Ranking normalizado de desempeño relativo; la infracción se determina separadamente.
&
Calidad, presión, continuidad, tratamiento, facturación y reclamos.
&
El ranking no genera por sí solo una sanción; el incumplimiento de estándares u obligaciones puede dar lugar a fiscalización y sanción separada.
&
Los planes de desarrollo y sus obras se fiscalizan fuera del ranking anual.
\\

Colombia
&
Estándares generales con metas anuales y gradualidades por prestador.
&
Cobertura, continuidad, calidad, satisfacción, tratamiento, planes y buenas prácticas.
&
Determinadas metas generan descuentos anuales en componentes tarifarios; otras consecuencias regulatorias operan por vías separadas.
&
El POIR controla activos planeados, ejecutados y puestos en operación.
\\
\hline
\end{tabularx}

\begin{minipage}{\textwidth}
\footnotesize
\textit{Nota:} La tabla resume los componentes predominantes y no reproduce la totalidad de las obligaciones regulatorias ni de los mecanismos de fiscalización y sanción aplicables en cada jurisdicción.
\end{minipage}
\end{table*}

\subsection{Síntesis comparativa y correspondencia con el caso peruano}
\label{anexo:tratamiento_indicadores_peru}
\label{anexo:sintesis_experiencia_internacional}

La revisión muestra que la selección del indicador y la fijación de su valor exigible son decisiones distintas. Colombia utiliza estándares generales con gradualidades; Italia asigna objetivos según clases; Ofwat define niveles dentro de cada revisión tarifaria; Portugal aplica una evaluación multidimensional basada en valores de referencia y calificaciones por entidad; y Chile utiliza un ranking normalizado de desempeño relativo. La experiencia no respalda una única regla de fijación. En los casos examinados, la referencia de desempeño se define considerando, con distinto grado de intensidad, las condiciones iniciales, la atribuibilidad del resultado y la información disponible.

Las consecuencias asociadas al desempeño tampoco son homogéneas. Italia y Colombia, junto con Inglaterra y Gales, incorporan mecanismos económicos vinculados con el incumplimiento de determinados resultados o compromisos; Chile separa el ranking comparativo de la determinación de infracciones; y Portugal utiliza principalmente la evaluación pública y comparativa, sin que una calificación desfavorable constituya por sí sola una sanción. La regularidad relevante no es, por tanto, el carácter reputacional de los indicadores, sino la vinculación de las consecuencias económicas o sancionadoras con dimensiones, compromisos u obligaciones identificables.

Pese a esa diversidad, los cinco marcos mantienen alguna separación entre la evaluación de las condiciones observadas y el control de las inversiones o actividades. Los resultados se miden mediante macroindicadores, compromisos, calificaciones, rankings o metas anuales; las inversiones quedan sujetas a programas, planes, entregables, cronogramas, puesta en operación y mecanismos propios de ajuste, devolución o sanción. A su vez, las dimensiones evaluadas conservan suficiente identificación, lo que limita la posibilidad de que la agregación oculte por completo resultados heterogéneos o de que el cumplimiento de una dimensión neutralice automáticamente la consecuencia asociada al incumplimiento de otra.

Esta distinción permite ordenar los principales indicadores utilizados en el Perú. Continuidad, presión, calidad del agua, tratamiento efectivo, agua no facturada y micromedición efectiva presentan equivalentes funcionales en los sistemas examinados. En cambio, los catastros, la instalación de equipos, el avance financiero y la ejecución de reservas se aproximan a capacidades, productos físicos o utilización de recursos, y requieren instrumentos específicos de seguimiento y fiscalización.

\begin{table*}[!t]
\centering
\caption{Correspondencia entre indicadores peruanos y el tratamiento observado internacionalmente}
\label{tab:anexo_correspondencia_peru_internacional}
\small
\renewcommand{\arraystretch}{1.16}
\setlength{\tabcolsep}{5pt}
\begin{tabularx}{\textwidth}{
>{\raggedright\arraybackslash}p{3.0cm}
>{\raggedright\arraybackslash}p{3.0cm}
>{\raggedright\arraybackslash}X
>{\raggedright\arraybackslash}p{3.35cm}}
\hline
\textbf{Variable en el Perú} &
\textbf{Objeto medido} &
\textbf{Tratamiento internacional comparable} &
\textbf{Implicancia para el ICG} \\
\hline

Continuidad, presión, calidad del agua y tratamiento efectivo
&
Resultado de la prestación
&
Indicadores centrales de continuidad, calidad técnica y cumplimiento sanitario o ambiental. Cuando existen consecuencias económicas o regulatorias, estas permanecen vinculadas a resultados u obligaciones identificables.
&
Incorporación directa, con trazabilidad por dimensión y localidad cuando corresponda.
\\

Agua no facturada y micromedición efectiva
&
Resultado operativo o intermedio
&
Pérdidas, fugas, consumo medido y confiabilidad de lecturas se utilizan como condiciones de eficiencia o sostenibilidad.
&
Incorporación posible cuando miden una condición observable y no la actividad ejecutada.
\\

Catastro técnico y comercial
&
Capacidad e instrumento de gestión
&
Los registros de activos y usuarios sustentan la planificación, el reporte y el aseguramiento de datos.
&
Seguimiento mediante hitos, cobertura, actualización, calidad y uso efectivo, fuera del índice principal.
\\

Instalación, renovación o reemplazo de medidores
&
Producto físico
&
La instalación puede verificarse como entregable; el resultado se aproxima mejor mediante conexiones medidas y lecturas válidas.
&
Separar la obligación física de la meta de micromedición efectiva.
\\

Avance financiero del programa de inversiones
&
Utilización de recursos
&
El gasto se supervisa, pero la protección regulatoria utiliza principalmente entregables, hitos, alcance y entrada en operación.
&
Mantenerlo como señal de seguimiento y alerta, sin incidencia directa sobre el desempeño.
\\

Ejecución de reservas
&
Cumplimiento financiero
&
Los recursos afectados se controlan separadamente de las intervenciones y de los resultados ambientales o de resiliencia.
&
Diferenciar financiamiento, actividad ejecutada y resultado obtenido.
\\
\hline
\end{tabularx}

\begin{minipage}{\textwidth}
\footnotesize
\textit{Nota:} La clasificación corresponde al objeto predominantemente medido. Una intervención puede contribuir posteriormente a varios resultados sin que su ejecución física o financiera se convierta, por esa razón, en un resultado del servicio.
\end{minipage}
\end{table*}

En conjunto, la comparación respalda una arquitectura en la que el índice principal concentre resultados de desempeño, mientras que las inversiones, productos físicos, capacidades de gestión y variables financieras permanezcan sujetos a instrumentos propios de planificación, verificación y control. Esta separación no supone que los resultados deban tener un tratamiento exclusivamente reputacional: en los sistemas examinados, su incumplimiento puede generar ajustes económicos o activar procedimientos específicos, y la consecuencia conserva una vinculación identificable con la dimensión, el compromiso o la obligación correspondiente. La evidencia comparada aporta así un referente institucional para evaluar el diseño peruano, no una demostración causal de los efectos de la reforma.
